% GENERATED FILE. Edit canonical lessons, then run npm run generate:books.
% canonical-source-hash: fnv1a64:112229d6991c0892
% canonical-lessons: BN-C333-jukti, BN-C333-boktobyo, BN-C333-pokkho, BN-C333-biruddhe, BN-C333-somorthon-kora

\chapter{Making a Case}
\label{ch:bn-making-a-case}
\hlchaptermodality{333}

\begin{chapteropening}
\textbf{By the end of this chapter:} \emph{I can argue for a view and say which side I am on.}

Complete the last lesson of chapter 333: I can argue for a view and say which side I am on.
\end{chapteropening}

\section[\texorpdfstring{\bn{যুক্তি} (jukti)}{যুক্তি (jukti)}]{\bn{যুক্তি} (jukti) — an argument; reasoning}

\label{lesson:BN-C333-jukti}



\begin{quote}
Before the new one: say the Bengali for mind you, then the Bengali for to hold good.
\end{quote}

\subsection*{You'll want to know: \bn{যুক্তি}}
\textbf{\bn{যুক্তি}} — \emph{jukti} — \textquotedblleft{}an argument; reasoning\textquotedblright{}.

\begin{quote}
\bn{তোমার} \bn{যুক্তিটা} \bn{ভালো}, \bn{কিন্তু} \bn{এটা} \bn{সব} \bn{জায়গায়} \bn{খাটে} \bn{না।}
\end{quote}

\begin{quote}
\emph{tomār juktiṭā bhālo, kintu eṭā sôb jāygāy khāṭe nā.}
\end{quote}

\begin{quote}
Your argument is good, but it doesn't hold everywhere.
\end{quote}

\begin{cousinweb}
\textbf{\bn{যুক্তি}} (\emph{jukti}) is Sanskrit \emph{yukti}, a joining and so a reasoning that holds together (Monier-Williams). It is a reasoned case, not a quarrel: \bn{তর্ক} \bn{করা} is to argue, \bn{যুক্তি} is what you argue with.
\end{cousinweb}

\begin{grammarlens}[title={What you've built}]
One more word for giving reasons and opinions.
\end{grammarlens}

\subsection*{Your turn}
\begin{itemize}
\raggedright
  \item \emph{Say it:} \emph{jukti}
  \item \emph{Say it:} \emph{jukti}, once more
  \item \emph{Say it:} \emph{tomār juktiṭā bhālo, kintu eṭā sôb jāygāy khāṭe nā.} — Your argument is good, but it doesn't hold everywhere.
  \item \emph{Recall:} say the Bengali for mind you, then the Bengali for to hold good, then say \emph{jukti} again
\end{itemize}

\begin{tcolorbox}[breakable,colback=teal!4,colframe=teal!35!black,title={Before you move on}]
What does \bn{যুক্তি} mean? (\textquotedblleft{}An argument; reasoning\textquotedblright{}.) Say it once more.
\end{tcolorbox}
\section[\texorpdfstring{\bn{বক্তব্য} (bôktôbyô)}{বক্তব্য (bôktôbyô)}]{\bn{বক্তব্য} (bôktôbyô) — what one has to say; a statement}

\label{lesson:BN-C333-boktobyo}



\begin{quote}
Before the new one: say the Bengali for to hold good, then the Bengali for an argument.
\end{quote}

\subsection*{You'll want to know: \bn{বক্তব্য}}
\textbf{\bn{বক্তব্য}} — \emph{bôktôbyô} — \textquotedblleft{}what one has to say; a statement\textquotedblright{}.

\begin{quote}
\bn{তোমার} \bn{বক্তব্য} \bn{আমরা} \bn{সবাই} \bn{শুনলাম।}
\end{quote}

\begin{quote}
\emph{tomār bôktôbyô āmrā sôbāi shunlām.}
\end{quote}

\begin{quote}
We all heard what you had to say.
\end{quote}

\begin{cousinweb}
\textbf{\bn{বক্তব্য}} (\emph{bôktôbyô}) is Sanskrit \emph{vaktavya}, 'what is to be said' (Monier-Williams). In Bengali it is the point a person puts forward in a meeting or a debate.
\end{cousinweb}

\begin{grammarlens}[title={What you've built}]
One more word for giving reasons and opinions.
\end{grammarlens}

\subsection*{Your turn}
\begin{itemize}
\raggedright
  \item \emph{Say it:} \emph{bôktôbyô}
  \item \emph{Say it:} \emph{bôktôbyô}, once more
  \item \emph{Say it:} \emph{tomār bôktôbyô āmrā sôbāi shunlām.} — We all heard what you had to say.
  \item \emph{Recall:} say the Bengali for to hold good, then the Bengali for an argument, then say \emph{bôktôbyô} again
\end{itemize}

\begin{tcolorbox}[breakable,colback=teal!4,colframe=teal!35!black,title={Before you move on}]
What does \bn{বক্তব্য} mean? (\textquotedblleft{}What one has to say; a statement\textquotedblright{}.) Say it once more.
\end{tcolorbox}
\section[\texorpdfstring{\bn{পক্ষ} (pôkkhô)}{পক্ষ (pôkkhô)}]{\bn{পক্ষ} (pôkkhô) — a side in a dispute; a party}

\label{lesson:BN-C333-pokkho}



\begin{quote}
Before the new one: say the Bengali for an argument, then the Bengali for a statement.
\end{quote}

\subsection*{You'll want to know: \bn{পক্ষ}}
\textbf{\bn{পক্ষ}} — \emph{pôkkhô} — \textquotedblleft{}a side in a dispute; a party\textquotedblright{}.

\begin{quote}
\bn{আমি} \bn{তোমার} \bn{পক্ষে} \bn{আছি।}
\end{quote}

\begin{quote}
\emph{āmi tomār pôkkhe āchhi.}
\end{quote}

\begin{quote}
I'm on your side.
\end{quote}

\begin{cousinweb}
\textbf{\bn{পক্ষ}} (\emph{pôkkhô}) is Sanskrit \emph{pakṣa}, a wing, and so a side or party (Monier-Williams). To be on someone's side is -\bn{এর} \bn{পক্ষে} \bn{থাকা}.
\end{cousinweb}

\begin{grammarlens}[title={What you've built}]
One more word for giving reasons and opinions.
\end{grammarlens}

\subsection*{Your turn}
\begin{itemize}
\raggedright
  \item \emph{Say it:} \emph{pôkkhô}
  \item \emph{Say it:} \emph{pôkkhô}, once more
  \item \emph{Say it:} \emph{āmi tomār pôkkhe āchhi.} — I'm on your side.
  \item \emph{Recall:} say the Bengali for an argument, then the Bengali for a statement, then say \emph{pôkkhô} again
\end{itemize}

\begin{tcolorbox}[breakable,colback=teal!4,colframe=teal!35!black,title={Before you move on}]
What does \bn{পক্ষ} mean? (\textquotedblleft{}A side in a dispute; a party\textquotedblright{}.) Say it once more.
\end{tcolorbox}
\section[\texorpdfstring{\bn{বিরুদ্ধে} (biruddhe)}{বিরুদ্ধে (biruddhe)}]{\bn{বিরুদ্ধে} (biruddhe) — against}

\label{lesson:BN-C333-biruddhe}



\begin{quote}
Before the new one: say the Bengali for a statement, then the Bengali for a side in a dispute.
\end{quote}

\subsection*{You'll want to know: \bn{বিরুদ্ধে}}
\textbf{\bn{বিরুদ্ধে}} — \emph{biruddhe} — \textquotedblleft{}against\textquotedblright{}.

\begin{quote}
\bn{গ্রামের} \bn{সবাই} \bn{নতুন} \bn{আইনের} \bn{বিরুদ্ধে।}
\end{quote}

\begin{quote}
\emph{grāmer sôbāi nôtun āiner biruddhe.}
\end{quote}

\begin{quote}
Everyone in the village is against the new law.
\end{quote}

\begin{cousinweb}
\textbf{\bn{বিরুদ্ধে}} (\emph{biruddhe}) is Sanskrit \emph{viruddha}, opposed or hindered (Monier-Williams), with the Bengali ending -\bn{এ}. It follows a word with -\bn{এর}: \bn{আইনের} \bn{বিরুদ্ধে}, against the law.
\end{cousinweb}

\begin{grammarlens}[title={What you've built}]
One more word for giving reasons and opinions.
\end{grammarlens}

\subsection*{Your turn}
\begin{itemize}
\raggedright
  \item \emph{Say it:} \emph{biruddhe}
  \item \emph{Say it:} \emph{biruddhe}, once more
  \item \emph{Say it:} \emph{grāmer sôbāi nôtun āiner biruddhe.} — Everyone in the village is against the new law.
  \item \emph{Recall:} say the Bengali for a statement, then the Bengali for a side in a dispute, then say \emph{biruddhe} again
\end{itemize}

\begin{tcolorbox}[breakable,colback=teal!4,colframe=teal!35!black,title={Before you move on}]
What does \bn{বিরুদ্ধে} mean? (\textquotedblleft{}Against\textquotedblright{}.) Say it once more.
\end{tcolorbox}
\section[\texorpdfstring{\bn{সমর্থন} \bn{করা} (sômôrthôn kôrā)}{সমর্থন করা (sômôrthôn kôrā)}]{\bn{সমর্থন} \bn{করা} (sômôrthôn kôrā) — to support (a view or a person); to back}

\label{lesson:BN-C333-somorthon-kora}



\begin{quote}
Before the new one: say the Bengali for a side in a dispute, then the Bengali for against.
\end{quote}

\subsection*{You'll want to know: \bn{সমর্থন} \bn{করা}}
\textbf{\bn{সমর্থন} \bn{করা}} — \emph{sômôrthôn kôrā} — \textquotedblleft{}to support (a view or a person); to back\textquotedblright{}.

\begin{quote}
\bn{আমরা} \bn{সবাই} \bn{মীরাকে} \bn{সমর্থন} \bn{করলাম।}
\end{quote}

\begin{quote}
\emph{āmrā sôbāi mīrāke sômôrthôn kôrlām.}
\end{quote}

\begin{quote}
We all backed Mira.
\end{quote}

Past (he/she …): \bn{তিনি} \bn{আমাদের} \bn{যুক্তি} \bn{সমর্থন} \bn{করলেন}, \emph{tini āmāder jukti sômôrthôn kôrlen}.

\begin{cousinweb}
\textbf{\bn{সমর্থন} \bn{করা}} (\emph{sômôrthôn kôrā}) is Sanskrit \emph{samarthana}, which Monier-Williams gives both as deliberation and as establishing or vindicating a claim, with \bn{করা} 'to do'. Bengali uses it for standing behind a view or a person.
\end{cousinweb}

\begin{grammarlens}[title={What you've built}]
One more word for giving reasons and opinions.
\end{grammarlens}

\subsection*{Your turn}
\begin{itemize}
\raggedright
  \item \emph{Say it:} \emph{sômôrthôn kôrā}
  \item \emph{Say it:} \emph{sômôrthôn kôrā}, once more
  \item \emph{Say it:} \emph{āmrā sôbāi mīrāke sômôrthôn kôrlām.} — We all backed Mira.
  \item \emph{Recall:} say the Bengali for a side in a dispute, then the Bengali for against, then say \emph{sômôrthôn kôrā} again
\end{itemize}

\begin{tcolorbox}[breakable,colback=teal!4,colframe=teal!35!black,title={Before you move on}]
What does \bn{সমর্থন} \bn{করা} mean? (\textquotedblleft{}To support (a view or a person); to back\textquotedblright{}.) Say it once more.
\end{tcolorbox}
